Subspace-major codes align neighbouring lanes/candidates; splitting reduces LUT storage.

\subsection{Address mapping and four-byte packing}

For list $\ell$, row $r$ maps to
\[
 p=o_\ell+r,\qquad \operatorname{offset}(Z[m,p])=mN+p.
\]
At fixed $m$, $\Delta\operatorname{offset}=\Delta p$. List boundaries can break adjacency.

Pack four subspace bytes per candidate:
\begin{equation}
 \mathsf{w}[j,p]=\sum_{t=0}^{3}2^{8t}Z[4j+t,p],\qquad
 \operatorname{offset}_{\rm bytes}(\mathsf{w}[j,p])=4(jN+p).
 \label{eq:5-packed-layout}
\end{equation}
For $0\leq t<4$ and eight-bit $Z$,
\[
 (\mathsf w[j,p]\gg8t)\mathbin{\&}255=Z[4j+t,p].
\]
Nibbles index \cref{eq:5-lut-factorisation}.

\par\medskip\noindent\begin{minipage}{\linewidth}
\captionsetup{hypcap=false}
\centering
% Address schematic; four lanes stand for a subset of a warp.
\begin{tikzpicture}[font=\small,
 box/.style={draw=ThesisNavy!65,rounded corners=2pt,align=center,
   text width=2.6cm,minimum height=1.2cm,inner sep=4pt},
 arr/.style={-{Latex[length=2mm]},draw=ThesisNavy!80,thick}]
\foreach \a in {0,1,2,3} {
 \node[box,fill=ThesisNavy!5] (w\a) at (3.3*\a,0)
  {lane $\a$\\$\mathsf{w}[j,p+\a]$};
}
\node[align=center] at (4.95,1.05)
 {Fixed $j$: adjacent 32-bit words across candidate positions};
\node[box,text width=11.5cm,minimum height=1.6cm] (decode) at (4.95,-2.5)
 {Unpack each lane's word independently at its position $p_a$\\
  $Z[4j,p_a]\;|\;Z[4j+1,p_a]\;|\;Z[4j+2,p_a]\;|\;Z[4j+3,p_a]$\\
  Four code bytes $\longrightarrow$ shifts and masks $\longrightarrow$ LUT indices};
\draw[arr] (w0.south) -- (decode.north west);
\draw[arr] (w1.south) -- (decode.north);
\draw[arr] (w2.south) -- (decode.north);
\draw[arr] (w3.south) -- (decode.north east);
\end{tikzpicture}

\captionof{figure}[Packed candidate-parallel access]{Packed candidate-parallel access, showing a subset of warp lanes.
At fixed word group $j$, consecutive positions access adjacent four-byte
words. Each word contains four subspace codes of one candidate. Unpacking
preserves each code byte; LUT splitting happens afterwards. The diagram
illustrates addresses, not a measured number of memory transactions.}
\label{fig:5-packed-access}
\end{minipage}\par\medskip

Construction maps $[M,N]$ bytes to $[M/4,N]$ words ($4\mid M$). Casting would combine candidate bytes. Code-byte costs:
\[
 R=W=NM,\qquad S_{\rm old+new}=2NM.
\]
Repacking is outside search timing.

\subsection{What packed access saves, and what it preserves}

One word load/unpack replaces four byte loads. The logical warp window is $32\times4=128$ bytes; transactions depend on alignment/active lanes. Payload stays $M$ bytes per candidate. Both arms are subspace-major, so transpose/cache effects are not isolated.

